\documentclass{article}

\usepackage{nips15submit_e,times}
\usepackage{hyperref}
\usepackage{url}
\usepackage{natbib}

\title{Literature Review: Multimodal Spatiotemporal\\Prediction of Peatland Fire Hotspots in Riau}

\author{
Author \\
\texttt{email@example.com}
}

\begin{document}

\maketitle

\begin{abstract}
This literature review surveys prior work on forest and peatland fire
prediction using spatiotemporal deep learning and multimodal fusion. We
identify three primary research gaps: (1) no systematic comparison between
tabular and spatiotemporal approaches exists for peatland hotspot prediction
in Indonesia, (2) multimodal fusion integrating satellite radar, weather,
rainfall, land cover, and peat vulnerability maps has not been explored in a
spatiotemporal framework, and (3) analysis of the relative contribution of
environmental factors to fire risk using explainable AI remains limited for
the Indonesian context.
\end{abstract}

\section{Literature Review}

\subsection{Limitations of Optical and Thermal Satellite Detection}

Conventional peatland fire detection relies on polar-orbiting satellites
such as MODIS (Moderate Resolution Imaging Spectroradiometer) and VIIRS
(Visible Infrared Imaging Radiometer Suite), which detect thermal anomalies
on the Earth's surface. Although effective for large surface fires, this
approach has three major limitations for peatland contexts.

First, MODIS has a 1~km spatial resolution per pixel, too coarse to detect
small-scale \emph{smoldering} fires common in the early stages of peat
combustion~\cite{vanwees2022}. Second, persistent cloud cover in tropical
regions, especially during fire season, causes many detections to be missed
by optical satellites. Third, peat fires burning below the surface have
substantially lower surface temperatures than open flames, frequently falling
below the detection threshold of MODIS algorithms~\cite{whitburn2016}.
Temporal hotspot patterns in Riau have been studied by~\cite{sitanggang2018},
who identified correlations between hotspot sequences and peatland conditions,
but that study did not integrate multimodal data. These limitations
demonstrate that no single sensor modality can reliably predict peatland
fires, motivating the need for a multimodal fusion approach.

\subsection{Deep Learning Approaches for Fire Prediction}

Advances in deep learning have opened new approaches for fire prediction
beyond thermal anomaly detection. The Convolutional LSTM (ConvLSTM),
introduced by~\cite{shi2015} for precipitation nowcasting, has become a
foundational architecture for various spatiotemporal prediction tasks,
including wildfire prediction.

Several recent studies have demonstrated promising ConvLSTM performance for
fire prediction.~\cite{masrur2024} developed an attention-based ConvLSTM to
predict wildfire spread in the United States, showing that spatiotemporal
deep learning significantly outperforms conventional Auto-Regressive Moving
Average (ARIMA) models.~\cite{chakraa2026} accelerated fire spread prediction
using ConvLSTM for operational guidance, achieving accurate forecasts at a
fraction of the computational cost of physics-based simulations.

For the multimodal setting,~\cite{papakis2025} proposed a multimodal ensemble
deep learning framework combining satellite imagery, topographic data, and
meteorological variables for fire prediction in Greece, achieving
competitive accuracy against state-of-the-art methods.
~\cite{mamun2026} developed a unified spatiotemporal deep learning framework
integrating 3D CNN, LSTM, and MLP for probabilistic wildfire prediction in
the Mediterranean region, attaining an AUC-ROC of 0.962.
~\cite{sinato2026} developed a fire risk prediction system in Argentina using
Dynamic World as a land cover source and FIRMS as validation labels, and won
the NASA Space Apps Challenge 2025.

\subsection{Multimodal Fusion for Environmental Prediction}

Multimodal fusion has become increasingly popular in environmental prediction
because it captures complementary signals from multiple physical sources. For
peatland fires, spatiotemporal fusion of radar, weather, rainfall, and land
cover data remains largely unexplored in the literature. In particular, the
integration of cloud-penetrating radar for direct soil moisture measurement
alongside continuous atmospheric reanalysis and land cover information has
not been evaluated for Indonesia's peatlands~\cite{taufik2023}.

\subsection{Research Gaps}

Based on the literature surveyed above, several research gaps can be
identified. First, no study has systematically compared tabular and
spatiotemporal approaches for peatland hotspot prediction in Indonesia.
Second, multimodal fusion integrating radar, weather, rainfall, land cover,
and peat vulnerability maps has not been conducted in a spatiotemporal
framework. Third, the relative contribution of environmental factors to
peatland fire risk in Indonesia, analyzed through explainable AI methods,
remains underexplored.

\bibliographystyle{unsrt}
\bibliography{references}

\end{document}
